\section*{Hoofdstuk \ref{ch:b2:geological-hazards} --- Geologische gevaren en risico}

\begin{solution}{exo:b2:geological-hazards:1}
(a) $1 - e^{-30/100} = 0.259$. (b) $e^{-100/100} = e^{-1} = 0.368$: meer dan één eeuw op drie
verstrijkt zonder. (c) $T = -50/\ln 0.95 = 975$ jaar.
\end{solution}

\begin{solution}{exo:b2:geological-hazards:2}
(a) Stad A: $(1/200) \times 5000 \times 0.02 = 0.5$ gebouw per jaar; stad B: $(1/50) \times 800
\times 0.10 = 1.6$. (b) Stad B, met 6.25 keer minder blootgestelde gebouwen. (c) Het gevaar is er 4 keer
frequenter en de kwetsbaarheid 5 keer hoger: samen een factor 20, die zwaarder weegt dan de
kleinere \omterm{def:b2:geological-hazards:risk}{blootstelling}.
\end{solution}

\begin{solution}{exo:b2:geological-hazards:3}
(a) $\sqrt{9.81 \times 4000} = \qty{198}{m.s^{-1}}$; $\sqrt{9.81 \times 200} = \qty{44.3}{m.s^{-1}}$;
$\sqrt{9.81 \times 10} = \qty{9.90}{m.s^{-1}}$. (b) $\num{1.5e6}/198.1 = \qty{7570}{s}$, \qty{2.10}{h}.
(c) $800/3.6 = \qty{222.2}{m.s^{-1}}$, en $h = c^2/g = 222.2^2/9.81 = \qty{5030}{m}$.
\end{solution}

\begin{solution}{exo:b2:geological-hazards:4}
(a) $\tan\ang{33}/\tan\ang{25} = 0.6494/0.4663 = 1.39$. (b) $F = 1$ wanneer $\beta = \varphi' =
\ang{33}$. (c) Met $c' = 0$ is $F = (\cos^2\beta - r_u)\tan\varphi'/(\sin\beta\cos\beta)$, en $\cos^2\beta
- r_u = \cos^2\beta\,(1 - \gamma_w/\gamma)$, dus $F = (1 - \gamma_w/\gamma)\tan\varphi'/\tan\beta$:
het water halveert de wrijving. $(1 - 9.81/20) \times 1.393 = 0.510 \times 1.393 = 0.710$: de
helling bezwijkt.
\end{solution}

\begin{solution}{exo:b2:geological-hazards:5}
(a) $696/50 = 13.9$ per jaar. (b) $e^{-30/26.2} = 0.318$, tegenover 0.304 waargenomen: dichtbij. De
waargenomen waarde is iets lager omdat de geclusterde naschokken en dubbelbevingen korte
intervallen toevoegen. (c) $1 - e^{-7/26.2} = 0.234$.
\end{solution}

\begin{solution}{exo:b2:geological-hazards:6}
(a) De magnitudeterm is $-0.105 + 0.229 \times 0.5 = 0.0095$. Voor $d = 0$, 10, 30 en
\qty{100}{km} is $r = 5.57$, 11.45, 30.51 en \qty{100.2}{km}, $\log_{10}\mathrm{PGA} = -0.571$,
$-0.814$, $-1.145$ en $-1.547$, en de mediane PGA 0.269, 0.153, 0.0715 en \qty{0.0284}{g}.
(b) $0.153 \times 10^{0.251} = 0.153 \times 1.78 = \qty{0.272}{g}$. (c) Eén standaardafwijking
aan weerszijden: $0.153/1.70 = \qty{0.0900}{g}$ tot $0.153 \times 1.70 = \qty{0.260}{g}$.
\end{solution}

\begin{solution}{exo:b2:geological-hazards:7}
(a) Gerangschikt: 420 ($T = 13/1 = 13$ jaar), 305 (6.5), 268 (4.33), 251 (3.25). (b) $305 + (420 -
305)\ln(10/6.5)/\ln(13/6.5) = 305 + 115 \times 0.621 = \qty{376}{m^3.s^{-1}}$. (c) De grootste piek
wordt uitgezet bij 13 jaar; elke waarde bij 100 jaar extrapoleert ver buiten de reeks, met een
extreemwaardenverdeling waarvan de 12 punten de keuze en de parameters nauwelijks vastleggen.
\end{solution}

\begin{solution}{exo:b2:geological-hazards:8}
(a) $0.6 \times 0.4 \times 0.5 \times 0.3 = 0.036$. (b) $P(\text{uitbarsting}) = 0.6 \times 0.4 = 0.24$;
zonder stroom in het dal $0.24 - 0.036 = 0.204$. (c) $0.6 \times 0.7 \times 0.5 \times 0.3 =
0.063$.
\end{solution}

\begin{solution}{exo:b2:geological-hazards:9}
Hier is $\sin\beta\cos\beta = 0.4698$, $\cos^2\beta = 0.6710$ en $\tan\ang{30}/\tan\ang{35} =
0.8245$. (a) $F = c'/(\gamma z\sin\beta\cos\beta) + 0.8245 = 1$ geeft $z_c = 8/(19 \times 0.4698
\times 0.1755) = \qty{5.11}{m}$. (b) $8/(19 \times 2 \times 0.4698) + 0.8245 = 0.448 + 0.825 =
1.27$. (c) $r_u = 9.81 \times 0.6710/19 = 0.3465$; $F = [8 + 38 \times (0.6710 - 0.3465) \times
0.5774]/(38 \times 0.4698) = 15.12/17.85 = 0.847$: de verzadigde laag bezwijkt.
\end{solution}

\begin{solution}{exo:b2:geological-hazards:10}
(a) $\Pr(\text{geen in } [0, s+t] \mid \text{geen in } [0, s]) = e^{-(s+t)/T}/e^{-s/T} = e^{-t/T}$.
(b) De 400 rustige jaren veranderen niets: $1 - e^{-50/475} = 0.0999$, \qty{10}{\percent}. (c) Ze
spreekt de aanname tegen dat de frequentie niet afhangt van de verstreken tijd: na de terugveer moet de
spanning opnieuw opbouwen, zodat de frequentie vlak na een aardbeving laag is en later groeit. Een
bronzone bevat veel breuken in verschillende stadia van hun cyclus; de som van veel onafhankelijke reeksen
gebeurtenissen gedraagt zich bijna als een Poisson-proces, en de frequentie is alles wat van elk bekend is.
\end{solution}

\begin{solution}{exo:b2:geological-hazards:11}
(a) $A = 0.5 \times (4000/10)^{1/4} = 0.5 \times 4.47 = \qty{2.24}{m}$. (b) $\lambda = cT$:
$198.1 \times 1200 = \qty{238}{km}$ in de diepe oceaan, $9.90 \times 1200 = \qty{11.9}{km}$ in
\qty{10}{m} water. (c) De amplitude is niet langer klein ten opzichte van de diepte (\qty{2.2}{m} in
\qty{10}{m}), zodat de golf steiler wordt en kan breken; een deel van de energie wordt teruggekaatst of gefocust door de
vorm van de zeebodem en van baaien; op het land hangt de \omterm{def:b2:geological-hazards:tsunami}{oploophoogte} af van de helling en de ruwheid van de
grond.
\end{solution}

\begin{solution}{exo:b2:geological-hazards:12}
(a) $10/5.8 + 4 = 1.72 + 4 = \qty{5.72}{s}$. (b) De S-golf bereikt \qty{50}{km} na $50/3.2 =
\qty{15.6}{s}$: $15.6 - 5.72 = \qty{9.90}{s}$ waarschuwingstijd; op \qty{100}{km} $31.3 - 5.72 =
\qty{25.5}{s}$. (c) $3.2 \times 5.72 = \qty{18.3}{km}$.
\end{solution}

\begin{solution}{pb:b2:geological-hazards:1}
Logaritmen worden in tussenstappen op vier decimalen gehouden.
\textbf{1.} $10^{4-5} = 0.1$, $10^{4-6} = 0.01$ en $10^{4-7} = 0.001$ per jaar.
\textbf{2.} 10, 100 en 1000 jaar.
\textbf{3.} $10^{4-7.5} = \num{3.16e-4}$ per jaar; $0.001 - \num{3.16e-4} = \num{6.84e-4}$ per jaar.
\textbf{4.} $0.1 - 0.01 = 0.09$ en $0.01 - 0.001 = 0.009$ per jaar.
\textbf{5.} $1 - e^{-0.01 \times 50} = 0.393$.
\textbf{6.} $1 - e^{-0.001 \times 50} = 0.0488$.
\textbf{7.} $r = \sqrt{10^2 + 5.57^2} = \qty{11.45}{km}$; $\log_{10} r = 1.0588$ en $-0.778 \times
1.0588 = -0.8237$.
\textbf{8.} $\mu = -0.105 + 0.229(M - 6) - 0.8237$: $-1.0432$, $-0.8142$ en $-0.6424$; medianen
0.0905, 0.153 en \qty{0.228}{g}.
\textbf{9.} $10^{-0.8142 + 0.230} = \qty{0.260}{g}$.
\textbf{10.} $z = (\log_{10} 0.3 + 0.8142)/0.230 = (-0.5229 + 0.8142)/0.230 = 1.266$, en $1 -
\Phi(1.266) = 0.103$.
\textbf{11.} $-0.105 + 0.229 \times 1.5 - 0.8237 = -0.5852$: \qty{0.260}{g}.
\textbf{12.} De kansen om \qty{0.2}{g} te overschrijden ($\log_{10} 0.2 = -0.6990$) zijn 0.0672
($z = 1.497$), 0.308 ($z = 0.501$) en 0.597 ($z = -0.246$): $\lambda = 0.09 \times 0.0672 + 0.009
\times 0.308 + \num{6.84e-4} \times 0.597 = \num{6.05e-3} + \num{2.77e-3} + \num{4.08e-4} =
\num{9.23e-3}$ per jaar.
\textbf{13.} Voor \qty{0.3}{g}: $z = 2.262$, 1.266 en 0.520, kansen 0.0118, 0.103 en
0.302; $\lambda = \num{1.07e-3} + \num{9.24e-4} + \num{2.06e-4} = \num{2.20e-3}$ per jaar.
\textbf{14.} Voor \qty{0.4}{g} ($\log_{10} 0.4 = -0.3979$): $z = 2.805$, 1.810 en 1.063,
kansen 0.00251, 0.0352 en 0.144; $\lambda = \num{2.26e-4} + \num{3.17e-4} + \num{9.84e-5} =
\num{6.41e-4}$ per jaar.
\textbf{15.} $-\ln 0.9/50 = \num{2.107e-3}$ per jaar, een \omterm{def:b2:geological-hazards:return}{terugkeerperiode} van 475 jaar.
\textbf{16.} Tussen \qty{0.3}{g} (\num{2.20e-3}) en \qty{0.4}{g} (\num{6.41e-4}).
\textbf{17.} De helling is $\ln(6.41/22.0)/\ln(4/3) = -1.233/0.2877 = -4.29$, en $a = 0.3 \times
(2.107/2.20)^{-1/4.29} = 0.3 \times 1.0101 = \qty{0.303}{g}$.
\textbf{18.} Bij \qty{0.3}{g} levert de schijf 5--6 \qty{48.5}{\percent} van de frequentie, 6--7
\qty{42.1}{\percent} en 7--7.5 slechts \qty{9.4}{\percent}. Zelfs de grootste aardbeving heeft een mediaan van
\qty{0.260}{g}, onder \qty{0.3}{g}: het niveau wordt alleen in de bovenste staart van de spreiding bereikt, en
de frequente middelgrote aardbevingen leveren er het grootste deel van.
\textbf{19.} $4/(19 \times 3 \times 0.4330) + \tan\ang{34}/\tan\ang{30} = 0.162 + 1.168 = 1.33$.
\textbf{20.} $r_u = 9.81 \times 0.75/19 = 0.387$; $F = [4 + 57 \times (0.75 - 0.387) \times
0.6745]/(57 \times 0.4330) = 17.95/24.68 = 0.727$: de verzadigde helling bezwijkt.
\textbf{21.} Met de grondwaterspiegel op een fractie $m$ van de dikte is $r_u = 0.387\,m$. $F = 1$ wanneer
$0.75 - 0.387\,m = (24.68 - 4)/(57 \times 0.6745) = 0.538$: $m = 0.548$, een grondwaterspiegel op
\qty{55}{\percent} van de laag.
\textbf{22.} Deling van $F = 1$ door $\gamma z$ geeft $0.0702 + (0.75 - 0.4330\,k) \times 0.6745 = 0.4330 +
0.75\,k$, dus $k_c = (0.0702 + 0.5059 - 0.4330)/(0.2921 + 0.75) = 0.1431/1.042 = 0.137$.
\textbf{23.} $\log_{10} 0.137 = -0.8633 = -0.105 + 0.229(M - 6) - 0.8237$ geeft $M = 6.29$. Zulke
aardbevingen komen voor met $10^{4 - 6.29} - 10^{-3.5} = \num{4.81e-3}$ per jaar, ongeveer eens in 208 jaar,
en de PGA van 475 jaar uit vraag~17 is meer dan het dubbele van de kritieke \qty{0.137}{g}: volgens deze
analyse bezwijkt de droge helling bij een schudden ver onder de beweging van 475 jaar.
\textbf{24.} De \omterm{def:b2:geological-hazards:psha}{piekgrondversnelling} met een kans van \qty{10}{\percent} om in 50 jaar te worden
overschreden in de stad is \textbf{\qty[text-series-to-math=true]{0.303}{g}}.
\end{solution}
